% Part: computability
% Chapter: recursive-functions
% Section: halting-problem

\documentclass[../../../include/open-logic-section]{subfiles}

\begin{document}

\olfileid{cmp}{rec}{hlt}
\olsection{నిలుపు సమస్య (హాల్టింగ్ సమస్య)}

సాధారణంగా \emph{నిలుపు సమస్య} అంటే, ఒక గణనీయ ప్రమేయపు
నిర్దేశం~$e$ (ఉదాహరణకు, ప్రోగ్రాము), ఒక సంఖ్య~$n$ ఇచ్చినప్పుడు,
ఇన్‌పుట్~$n$పై ఆ ప్రమేయపు గణన ఆగి ఫలితాన్ని ఇస్తుందో లేదో నిర్ణయించే
సమస్య. ప్రసిద్ధంగా, అలన్ ట్యూరింగ్ ఈ సమస్యనే ఒక గణనీయ ప్రమేయంతో
పరిష్కరించలేమని నిరూపించాడు. అంటే కింది ప్రమేయం గణనీయం కాదు:
\[
h(e, n) =
\begin{cases}
1 & \text{ఇన్‌పుట్ $n$పై గణన $e$ ఆగితే}\\
0 & \text{ఇతర సందర్భాల్లో.}
\end{cases}
\]

పాక్షిక పునరావృత్త ప్రమేయాల సందర్భంలో, ప్రోగ్రాము నిర్దేశపు పాత్రను
క్లీనీ నియత రూప సిద్ధాంతంలోని సూచిక~$e$ పోషించవచ్చు. $f$ పాక్షిక
పునరావృత్త ప్రమేయమైతే, నియత రూప సిద్ధాంతపు సమీకరణ నెరవేరే ఏ~$e$ అయినా
$f$కు ఒక సూచిక. ఒక సంఖ్య~$e$ ఇచ్చినప్పుడు,
\[
\cfind{e}(x) \simeq U(\mu s \; T(e, x, s))
\]
పాక్షిక పునరావృత్త ప్రమేయమని నియత రూప సిద్ధాంతం చెబుతుంది; అంతేకాక,
ప్రతి పాక్షిక పునరావృత్త ప్రమేయం $f\colon\Nat\to\Nat$కు, ప్రతి
$x\in\Nat$కూ $\cfind{e}(x)\simeq f(x)$ అయ్యే $e\in\Nat$ ఒకటి
ఉంటుంది. నిజానికి ప్రతి అటువంటి $f$కు ఒక్కటి మాత్రమే కాక అనంతంగా
అనేక $e$లు ఉంటాయి. \emph{నిలుపు ప్రమేయం}~$h$ను ఇలా నిర్వచిస్తాం:
\[
h(e, x) =
\begin{cases}
1 & \text{$\cfind{e}(x)\fdefined$ అయితే}\\
0 & \text{$\cfind{e}(x)\fundefined$ అయితే.}
\end{cases}
\]
అంటే $h(e,x)=0$ అయినప్పుడే $\cfind{e}(x)\fundefined$.

\begin{thm}
\ollabel{thm:halting-problem}
నిలుపు ప్రమేయం~$h$ పాక్షిక పునరావృత్త ప్రమేయం కాదు.
\end{thm}

\begin{proof}
$h$ పాక్షిక పునరావృత్త ప్రమేయమైతే, కింది ప్రమేయాన్ని నిర్వచించవచ్చు:
\[
d(y) =
\begin{cases}
1 & \text{$h(y, y) = 0$ అయితే}\\
\umin{x}{x \neq x} & \text{ఇతర సందర్భాల్లో.}
\end{cases}
\]
$x\neq x$ను నెరవేర్చే సంఖ్య $x$ ఏదీ లేదు కనుక
$\umin{x}{x\neq x}$ లేదు. అందువల్ల $h(y,y)\neq0$ అయినప్పుడే
$d(y)\fundefined$. ఈ నిర్వచనం నుంచి
\begin{enumerate}
\item $\cfind{y}(y)\fundefined$ అయినప్పుడే $d(y)\fdefined$;
\item $\cfind{y}(y)\fdefined$ అయినప్పుడే $d(y)\fundefined$
\end{enumerate}
అని వస్తుంది. $h$ పాక్షిక పునరావృత్త ప్రమేయమైతే $d$ కూడా అలాంటిదే.
కాబట్టి క్లీనీ నియత రూప సిద్ధాంతం ప్రకారం దానికి ఒక సూచిక~$e_d$
ఉంటుంది. $h(e_d,e_d)$ విలువను పరిశీలించండి. అది $0$ లేదా $1$ అనే
రెండు సందర్భాలే ఉన్నాయి.
\begin{enumerate}
\item $h(e_d,e_d)=1$ అయితే $\cfind{e_d}(e_d)\fdefined$. కానీ
  $\cfind{e_d}\simeq d$; మరోవైపు $h(e_d,e_d)\neq0$ కనుక
  $d(e_d)\fundefined$. ఇది వైరుధ్యం.
\item $h(e_d,e_d)=0$ అయితే $\cfind{e_d}(e_d)\fundefined$. కానీ
  అప్పుడు $d(e_d)\fdefined$; ఇది మళ్లీ $\cfind{e_d}\simeq d$కు
  వైరుధ్యం.
\end{enumerate}
అందువల్ల $h(e_d,e_d)$కు సాధ్యమైన రెండు విలువలూ వైరుధ్యాన్నే ఇస్తాయి.
$h$ పాక్షిక పునరావృత్త ప్రమేయం అనే ఊహ తప్పు; కనుక నిలుపు ప్రమేయం
$h$ పాక్షిక పునరావృత్త ప్రమేయం కాదు.
\sourcecorrection{OLTECRREM-010}{ఆంగ్ల మూలం ప్రతి $e$కు $\cfind{e}$ను పాక్షిక పునరావృత్త ప్రమేయంగా నిర్వచించిన వెంటనే ``$e$ అసలు సూచిక కాకపోవచ్చు'' అనే అదనపు సందర్భాన్ని చేర్చి, వికర్ణ నిరూపణ చివర దానిపైనే ఆధారపడింది. అదే మూలంలోని సార్వత్రిక $\cfind{e}$ కుటుంబానికి అనుగుణంగా ఆ అసంగత సందర్భాన్ని తొలగించి, $\cfind{e_d}\simeq d$తో రెండు విలువల నుంచీ నేరుగా వచ్చే వైరుధ్యాలను స్పష్టం చేశాం.}
\end{proof}

\end{document}
